As the longest shadow of the Sun on the given day is of finite length, the
Sun is circumpolar for this observer on this day. \hfill[2.0]

% FIGURE NEEDED: two-panel diagram -- left: stick OA with shadow OS and Sun
% altitude theta; right: celestial hemisphere with the Sun's circumpolar
% diurnal circle, angles phi, 90-delta_sun, and the two extreme altitudes
% theta_1, theta_2 (2016_theory_sol.pdf, page 4).
In the figure, the left panel shows the shadow $OS$ formed by stick $OA$ (of
length 1.000\,m), and the right panel shows the Sun's location in two cases.

For an altitude $\theta$ of the Sun,
\[ \tan\theta = \frac{OA}{OS} = \frac{1.000\ \mathrm{m}}{OS} \]
\[ \therefore\ \cot\theta = OS\ \text{(in metres)} \]
\hfill[1.0]

Let $\theta_1$ and $\theta_2$ be the altitude in the two extreme cases.
\[ \theta_1 = 180^\circ - \phi - (90^\circ - \delta_\odot)
   = 90^\circ - \phi + \delta_\odot \]
\hfill[1.0]
\begin{align*}
\cot(90^\circ - \phi + \delta_\odot) &= 1.732\\
\therefore\ \tan(\phi - \delta_\odot) &= 1.732\\
\phi - \delta_\odot &= \tan^{-1}(1.732) = 60^\circ = 1.047\ \mathrm{rad}
\end{align*}
\hfill[1.5]
\[ \theta_2 = \phi - (90^\circ - \delta_\odot)
   = \phi - 90^\circ + \delta_\odot \]
\hfill[1.0]
\begin{align*}
\cot(\phi - 90^\circ + \delta_\odot) &= 5.671\\
\therefore\ \tan(\phi - 90^\circ + \delta_\odot) &= \frac{1}{5.671}\\
\phi + \delta_\odot &= \tan^{-1}\left(\frac{1}{5.671}\right) + 90^\circ
  = 100^\circ = 1.745\ \mathrm{rad}
\end{align*}
\hfill[1.5]

Solving,
\[ \boxed{\phi = 80^\circ} = 1.396\ \mathrm{rad} \]
\hfill[1.0]
\[ \boxed{\delta_\odot = 20^\circ} = 0.349\ \mathrm{rad} \]
\hfill[1.0]

Given high accuracy of shadow length, only $\pm 0.5^\circ$ is allowed.

One can also solve the question by manipulating $\tan(\phi - \delta_\odot)$
and $\tan(\phi + \delta_\odot)$, to get $\tan(\phi)$ and $\tan(\delta_\odot)$.
